% !PRINTED 538
\weberpage{538}

Wir setzen $a$ als ungerade voraus und lassen $\varphi$ ein volles
Repräsentantensystem der Diskriminante $-m$ durchlaufen. Ist
dann $m \equiv 3 \pmod 8$, so ist $c$ ungerade, und $\varphi_1$, $\varphi_2$, $\varphi_3$
durchlaufen zusammen ein Repräsentantensystem der Diskriminante $-4m$. Denn unter den $3h$ Wurzeln dieser Formen
\[
\frac{\omega}{2}, \quad 2\,\omega, \quad 1 - \frac{2}{\omega+1}
\]
kommen nach § 123 keine äquivalenten vor.

Ist dagegen $m \equiv 7 \pmod 8$, so ist $c$ gerade, $\frac{1}{2}\varphi_1$, $\frac{1}{2}\varphi_3$ durchlaufen je ein Repräsentantensystem der Diskriminante $-m$, $\varphi_2$
durchläuft ein Repräsentantensystem der Diskriminante $-4m$.

Durchläuft also $\psi$ ein Repräsentantensystem der Diskriminante $-4m$, und setzen wir
\weq{13}{S = \overset{x,y}{\Sigma}\, \frac{1}{\psi^s},}
so ergibt sich aus (10):
\weq{14}{\begin{aligned}
\Bigl(1 + \frac{2}{4^s}\Bigr) \overset{\varphi}{\Sigma} S' &= \overset{\psi}{\Sigma} S, & m &\equiv 3 \\
\Bigl(1 + \frac{2}{4^s} - \frac{2}{2^s}\Bigr) \overset{\varphi}{\Sigma} S' &= \overset{\psi}{\Sigma} S, & m &\equiv 7
\end{aligned} \quad (\mathrm{mod}\ 8).}
Setzen wir in den Formeln (29), (30) des vorigen Paragraphen
\[
\psi' = 2\,\varphi, \quad \psi = (A, 2B, C),
\]
also:
\[
\begin{aligned}
A &= a, \\
B &= b - a, \\
C &= 4c - 2b + a,
\end{aligned}
\]
so ergibt sich:
\[
S - \frac{1}{2^s} S' = \frac{2\pi}{\sqrt{m}} \log \frac{f(\omega_1)f(\omega_2)}{\sqrt{2}},
\]
und wenn $\psi$ ein Repräsentantensystem der Diskriminante $-4m$
durchläuft, so durchläuft $\varphi$ dreimal oder nur einmal ein Repräsentantensystem der Diskriminante $-m$, je nachdem $m \equiv 3$
oder $\equiv 7 \pmod 8$ ist. Wir erhalten also, wenn $\varepsilon$ in diesen
beiden Fällen 3 oder 1 ist:
\weq{15}{\operatorname*{Lim}_{s=1} \Bigl(\overset{\psi}{\Sigma} S - \frac{\varepsilon}{2^s} \overset{\varphi}{\Sigma} S'\Bigr) = \frac{4\pi}{\sqrt{m}} \log \Pi \frac{f(\omega)}{\sqrt[4]{2}},}
worin $\omega$ die Wurzeln der Formen $\psi$ durchläuft. Und daraus
ergibt sich nach (14):
%CONT>>
